\documentclass[12pt,spanish]{article}
\usepackage[utf8]{inputenc}
\usepackage{amsmath, amssymb, amsthm}
\usepackage{geometry}
\usepackage[spanish]{babel}
\usepackage{enumitem}
\usepackage{xcolor}
\geometry{top=2.5cm, bottom=2.5cm, left=2.5cm, right=2.5cm}

\title{Manual Completo de Integrales Múltiples\\ (Dobles y Triples)}
\author{Universidad\\ Curso de Cálculo Multivariable}
\date{}

\newtheorem{definicion}{Definición}[section]
\newtheorem{teorema}{Teorema}[section]
\newtheorem{propiedad}{Propiedad}[section]
\newtheorem{ejemplo}{Ejemplo}[section]
\newtheorem{ejercicio}{Ejercicio}[section]
\newtheorem{solucion}{Solución}

\newcommand{\RR}{\mathbb{R}}
\newcommand{\iintD}{\iint_D}
\newcommand{\iiintD}{\iiint_D}
\newcommand{\dA}{\,dA}
\newcommand{\dV}{\,dV}

\begin{document}

\maketitle

\begin{abstract}
Este material presenta una exposición teórica profunda y una extensa colección de ejercicios sobre integrales dobles y triples. Se cubren los fundamentos, el teorema de Fubini, regiones de integración, cambio de orden, cambio de variable (incluyendo coordenadas polares, cilíndricas y esféricas), aplicaciones geométricas y físicas. Al final se incluye un ejercicio práctico completo con su respuesta.
\end{abstract}

\tableofcontents

\section{Integrales dobles: fundamentos}

\subsection{Definición y motivación}

La integral doble extiende el concepto de integral definida a funciones de dos variables. Geométricamente, si $f(x,y)\ge 0$, la integral $\iint_R f(x,y)\,dA$ representa el volumen del sólido limitado por la superficie $z=f(x,y)$ y la región plana $R$ en el plano $xy$.

\begin{definicion}
Sea $R=[a,b]\times[c,d]$ un rectángulo. Para una partición en subrectángulos $R_{ij}$ de área $\Delta A_{ij}=\Delta x_i\Delta y_j$, elegimos puntos muestra $(x_{ij}^*,y_{ij}^*)$ y formamos la suma de Riemann
\[
S=\sum_{i=1}^n\sum_{j=1}^m f(x_{ij}^*,y_{ij}^*)\Delta A_{ij}.
\]
Si el límite cuando la norma de la partición tiende a cero existe, dicho límite es la integral doble de $f$ sobre $R$:
\[
\iint_R f(x,y)\,dA = \lim_{\|P\|\to0} S.
\]
\end{definicion}

\subsection{Propiedades básicas}
\begin{propiedad}
La integral doble es lineal: $\iint_R (af+bg)=a\iint_R f+b\iint_R g$.
\end{propiedad}
\begin{propiedad}
Si $R=R_1\cup R_2$ con $R_1\cap R_2$ de medida nula, entonces $\iint_R f = \iint_{R_1}f+\iint_{R_2}f$.
\end{propiedad}
\begin{propiedad}
Si $f\le g$ en $R$, entonces $\iint_R f \le \iint_R g$.
\end{propiedad}

\subsection{Teorema de Fubini (evaluación mediante integrales iteradas)}
\begin{teorema}[Fubini]
Sea $f$ continua en $R=[a,b]\times[c,d]$. Entonces
\[
\iint_R f(x,y)\,dA = \int_a^b \left(\int_c^d f(x,y)\,dy\right)dx = \int_c^d \left(\int_a^b f(x,y)\,dx\right)dy.
\]
\end{teorema}

\subsection{Regiones de integración generales: tipo I y tipo II}

Una región $D$ puede no ser un rectángulo. Distinguimos dos tipos:

\begin{definicion}[Región tipo I]
\[
D = \{(x,y): a\le x\le b,\; g_1(x)\le y\le g_2(x)\},
\]
con $g_1,g_2$ continuas. La integral se escribe
\[
\iint_D f(x,y)\,dA = \int_a^b \int_{g_1(x)}^{g_2(x)} f(x,y)\,dy\,dx.
\]
\end{definicion}

\begin{definicion}[Región tipo II]
\[
D = \{(x,y): c\le y\le d,\; h_1(y)\le x\le h_2(y)\},
\]
con $h_1,h_2$ continuas. Entonces
\[
\iint_D f(x,y)\,dA = \int_c^d \int_{h_1(y)}^{h_2(y)} f(x,y)\,dx\,dy.
\]
\end{definicion}

\subsection{Cambio de orden de integración}
A veces conviene intercambiar el orden. Para ello se dibuja la región descrita por los límites originales y se reexpresa en el otro tipo.

\section{Cambio de variable en integrales dobles}

\subsection{Jacobiano}
Si $x=x(u,v), y=y(u,v)$ es una transformación biyectiva y diferenciable, entonces
\[
\iint_D f(x,y)\,dxdy = \iint_{D^*} f(x(u,v),y(u,v)) \left| \frac{\partial(x,y)}{\partial(u,v)} \right| du\,dv,
\]
donde el jacobiano es
\[
\frac{\partial(x,y)}{\partial(u,v)} = \det\begin{pmatrix}
\frac{\partial x}{\partial u} & \frac{\partial x}{\partial v} \\
\frac{\partial y}{\partial u} & \frac{\partial y}{\partial v}
\end{pmatrix}.
\]

\subsection{Coordenadas polares}
El cambio más importante:
\[
x = r\cos\theta,\quad y = r\sin\theta,\quad dA = r\,dr\,d\theta.
\]
Se usa cuando la región tiene simetría circular o el integrando contiene $x^2+y^2$.

\section{Integrales triples}

\subsection{Definición y extensión de Fubini}
Para una función $f(x,y,z)$ sobre una región $E\subset\RR^3$, la integral triple se define mediante sumas de Riemann. El teorema de Fubini permite evaluarla como integrales iteradas en cualquier orden, siempre que los límites se describan adecuadamente.

\subsection{Coordenadas cilíndricas}
\[
x = r\cos\theta,\; y = r\sin\theta,\; z = z,\quad dV = r\,dr\,d\theta\,dz.
\]

\subsection{Coordenadas esféricas}
\[
x = \rho\sin\phi\cos\theta,\; y = \rho\sin\phi\sin\theta,\; z = \rho\cos\phi,\quad dV = \rho^2\sin\phi\,d\rho\,d\phi\,d\theta,
\]
con $\rho\ge0$, $0\le\phi\le\pi$, $0\le\theta\le2\pi$.

\section{Aplicaciones de las integrales múltiples}

\subsection{Área y volumen}
\begin{itemize}
\item Área de una región plana $D$: $\displaystyle A(D)=\iint_D dA$.
\item Volumen bajo $z=f(x,y)\ge0$ sobre $D$: $\displaystyle V=\iint_D f(x,y)\,dA$.
\item Volumen entre dos superficies $z=f_{\text{sup}}(x,y)$ y $z=f_{\text{inf}}(x,y)$:
\[
V=\iint_D \bigl(f_{\text{sup}}(x,y)-f_{\text{inf}}(x,y)\bigr)\,dA.
\]
\item Volumen en el espacio: $\displaystyle V(E)=\iiint_E dV$.
\end{itemize}

\subsection{Masa y centro de masa}
Para una lámina plana con densidad $\delta(x,y)$:
\begin{align*}
\text{Masa } M &= \iint_D \delta(x,y)\,dA,\\
M_x &= \iint_D y\,\delta(x,y)\,dA,\quad M_y = \iint_D x\,\delta(x,y)\,dA,\\
\bar{x} &= M_y/M,\quad \bar{y}=M_x/M.
\end{align*}
Para un sólido con densidad $\delta(x,y,z)$:
\[
M=\iiint_E \delta\,dV,\quad
\bar{x}=\frac{1}{M}\iiint_E x\delta\,dV,\quad \text{etc.}
\]

\section{Ejercicios resueltos paso a paso}

\begin{ejercicio}
Calcular $\displaystyle\iint_R (x^2+y^2)\,dA$ donde $R=[0,1]\times[0,1]$.
\end{ejercicio}
\begin{solucion}
Usando Fubini:
\[
\iint_R (x^2+y^2)\,dxdy = \int_0^1\int_0^1 (x^2+y^2)\,dx\,dy = \int_0^1\left[\frac{x^3}{3}+xy^2\right]_{x=0}^{1}dy = \int_0^1\left(\frac13+y^2\right)dy = \left[\frac{y}{3}+\frac{y^3}{3}\right]_0^1 = \frac23.
\]
\end{solucion}

\begin{ejercicio}
Calcular $\displaystyle\iint_D xy\,dA$ donde $D$ es el triángulo con vértices $(0,0),(1,0),(0,1)$.
\end{ejercicio}
\begin{solucion}
$D$ es tipo I: $0\le x\le1$, $0\le y\le1-x$. Entonces
\[
\iint_D xy\,dA = \int_0^1\int_0^{1-x} xy\,dy\,dx = \int_0^1 x\left[\frac{y^2}{2}\right]_0^{1-x}dx = \frac12\int_0^1 x(1-x)^2dx.
\]
Expandiendo: $(1-x)^2=1-2x+x^2$, luego $x(1-x)^2=x-2x^2+x^3$. Integrando:
\[
\frac12\left[\frac{x^2}{2}-\frac{2x^3}{3}+\frac{x^4}{4}\right]_0^1 = \frac12\left(\frac12-\frac23+\frac14\right)=\frac12\left(\frac{6-8+3}{12}\right)=\frac12\cdot\frac1{12}=\frac1{24}.
\]
\end{solucion}

\begin{ejercicio}
Cambiar el orden de integración en $\displaystyle\int_0^1\int_{x^2}^{\sqrt{x}} f(x,y)\,dy\,dx$ y reescribir la integral.
\end{ejercicio}
\begin{solucion}
Los límites originales: $0\le x\le1$, $x^2\le y\le\sqrt{x}$. La región está entre $y=x^2$ e $y=\sqrt{x}$. Para $y$ fijo, $x$ va de $y^2$ a $\sqrt{y}$? Cuidado: despejando $x$ de $y=x^2$ da $x=\sqrt{y}$ (rama positiva) y de $y=\sqrt{x}$ da $x=y^2$. Como $y$ varía entre el mínimo y máximo de las curvas: cuando $x=0$, $y=0$; cuando $x=1$, $y=1$ y $y=1$. En la región $x^2\le y\le\sqrt{x}$, entonces $0\le y\le1$ y para cada $y$, $x$ está entre $y^2$ y $\sqrt{y}$ (pues $x^2\le y\Rightarrow x\le\sqrt{y}$ y $y\le\sqrt{x}\Rightarrow y^2\le x$). Luego
\[
\int_0^1\int_{x^2}^{\sqrt{x}} f(x,y)\,dy\,dx = \int_0^1\int_{y^2}^{\sqrt{y}} f(x,y)\,dx\,dy.
\]
\end{solucion}

\begin{ejercicio}
Calcular $\displaystyle\iint_D e^{-x^2-y^2}\,dA$ donde $D$ es el círculo $x^2+y^2\le a^2$.
\end{ejercicio}
\begin{solucion}
Usamos polares: $x=r\cos\theta$, $y=r\sin\theta$, $dA=r\,dr\,d\theta$. Los límites $0\le r\le a$, $0\le\theta\le2\pi$. Entonces
\[
\iint_D e^{-r^2}\,r\,dr\,d\theta = \int_0^{2\pi}\int_0^a r e^{-r^2}\,dr\,d\theta = 2\pi\cdot\left[-\frac12 e^{-r^2}\right]_0^a = 2\pi\cdot\frac12(1-e^{-a^2}) = \pi(1-e^{-a^2}).
\]
\end{solucion}

\begin{ejercicio}
Hallar el volumen del sólido acotado por el paraboloide $z=4-x^2-y^2$ y el plano $z=0$.
\end{ejercicio}
\begin{solucion}
El sólido está sobre el disco $x^2+y^2\le4$. En polares:
\[
V = \iint_D (4-x^2-y^2)\,dA = \int_0^{2\pi}\int_0^2 (4-r^2)\,r\,dr\,d\theta = 2\pi\int_0^2 (4r-r^3)dr = 2\pi\left[2r^2-\frac{r^4}{4}\right]_0^2 = 2\pi(8-4)=8\pi.
\]
\end{solucion}

\begin{ejercicio}
Calcular $\displaystyle\iiint_E z\,dV$ donde $E$ es el tetraedro con vértices $(0,0,0),(1,0,0),(0,1,0),(0,0,1)$.
\end{ejercicio}
\begin{solucion}
El tetraedro está limitado por $x\ge0$, $y\ge0$, $z\ge0$ y $x+y+z\le1$. Proyectando sobre $xy$: $0\le x\le1$, $0\le y\le1-x$, y $0\le z\le1-x-y$. Entonces
\[
\iiint_E z\,dV = \int_0^1\int_0^{1-x}\int_0^{1-x-y} z\,dz\,dy\,dx = \int_0^1\int_0^{1-x} \frac{(1-x-y)^2}{2}\,dy\,dx.
\]
Integramos respecto a $y$: sustitución $u=1-x-y$, $du=-dy$:
\[
\int_0^{1-x} \frac{(1-x-y)^2}{2}dy = \frac12\int_{u=1-x}^{0} u^2(-du) = \frac12\int_0^{1-x} u^2\,du = \frac12\cdot\frac{(1-x)^3}{3} = \frac{(1-x)^3}{6}.
\]
Luego
\[
\int_0^1 \frac{(1-x)^3}{6}dx = \frac16\cdot\frac14 = \frac1{24}.
\]
\end{solucion}

\begin{ejercicio}
Usar coordenadas cilíndricas para calcular $\displaystyle\iiint_E \sqrt{x^2+y^2}\,dV$, donde $E$ es la región dentro del cilindro $x^2+y^2=1$, entre $z=0$ y $z=1$.
\end{ejercicio}
\begin{solucion}
En cilíndricas: $x=r\cos\theta$, $y=r\sin\theta$, $z=z$, $dV=r\,dr\,d\theta\,dz$. La región es $0\le r\le1$, $0\le\theta\le2\pi$, $0\le z\le1$. El integrando $\sqrt{x^2+y^2}=r$. Entonces
\[
\iiint_E r\,dV = \int_0^{2\pi}\int_0^1\int_0^1 r\cdot r\,dz\,dr\,d\theta = \int_0^{2\pi}\int_0^1 r^2\,dr\,d\theta = 2\pi\cdot\frac13 = \frac{2\pi}{3}.
\]
\end{solucion}

\section{Ejercicios propuestos (para practicar)}

A continuación se presentan numerosos ejercicios sin resolver. Intente cada uno aplicando la teoría aprendida.

\begin{ejercicio}
Evalúe $\displaystyle\iint_R (3x+2y)\,dA$ donde $R=[0,2]\times[1,3]$.
\end{ejercicio}

\begin{ejercicio}
Calcule $\displaystyle\iint_D x^2y\,dA$ siendo $D$ la región acotada por $y=x^2$ e $y=2x$.
\end{ejercicio}

\begin{ejercicio}
Invierta el orden de integración en $\displaystyle\int_0^1\int_{2y}^2 e^{x^2}\,dx\,dy$ y evalúe.
\end{ejercicio}

\begin{ejercicio}
Encuentre el área de la región limitada por $y=x^2-2x$ e $y=x$.
\end{ejercicio}

\begin{ejercicio}
Calcule $\displaystyle\iint_D \frac{x}{1+y^2}\,dA$ donde $D$ es el triángulo de vértices $(0,0),(2,0),(0,1)$.
\end{ejercicio}

\begin{ejercicio}
Use coordenadas polares para evaluar $\displaystyle\iint_D \frac{1}{1+x^2+y^2}\,dA$, $D$: $x^2+y^2\le1$, $x\ge0$, $y\ge0$.
\end{ejercicio}

\begin{ejercicio}
Halle el volumen del sólido en el primer octante acotado por $z=1-x^2$, $y=0$, $y=2$, $z=0$.
\end{ejercicio}

\begin{ejercicio}
Calcule $\displaystyle\iiint_E (x^2+y^2)\,dV$ donde $E$ es la región limitada por $z=4-x^2-y^2$ y $z=1$.
\end{ejercicio}

\begin{ejercicio}
Use coordenadas esféricas para calcular $\displaystyle\iiint_E z\,dV$ donde $E$ es la semiesfera superior $x^2+y^2+z^2\le a^2$, $z\ge0$.
\end{ejercicio}

\begin{ejercicio}
Halle la masa de una lámina que ocupa la región $D$ acotada por $y=\sqrt{x}$, $y=0$, $x=4$ si la densidad es $\delta(x,y)=xy$.
\end{ejercicio}

\begin{ejercicio}
Determine el centro de masa de la lámina triangular con vértices $(0,0),(2,0),(0,2)$ y densidad constante $\delta=3$.
\end{ejercicio}

\begin{ejercicio}
Calcule el volumen del sólido que está dentro del cilindro $x^2+y^2=1$ y entre los planos $z=0$ y $z=x+2$.
\end{ejercicio}

\begin{ejercicio}
Evalúe $\displaystyle\iint_D \frac{x^2+y^2}{x^2}\,dxdy$ donde $D$ es la región del primer cuadrante entre $xy=1$, $xy=2$, $y=x$, $y=4x$ (sugerencia: cambio $u=xy$, $v=y/x$).
\end{ejercicio}

\begin{ejercicio}
Calcule la integral triple $\displaystyle\iiint_E 2x\,dV$ donde $E$ es la región acotada por $z=1-x^2$, $z=0$, $y=0$, $y=x+1$ en el primer octante.
\end{ejercicio}

\begin{ejercicio}
Use un cambio de variable adecuado para calcular $\displaystyle\iint_R (x+y)e^{x-y}\,dxdy$ donde $R$ es el paralelogramo con vértices $(0,0),(1,1),(2,0),(1,-1)$.
\end{ejercicio}

\section{Ejercicio práctico completo con respuesta}

A continuación se presenta un problema que integra varias técnicas. Resuélvalo paso a paso y luego verifique la respuesta.

\begin{ejercicio}[Práctico]
Sea $D$ la región del primer octante limitada por el cilindro $x^2+z^2=1$, el cono $z^2=x^2+y^2$ y el plano $y=0$.
\begin{enumerate}[label=\alph*)]
\item Describa la región $D$ en coordenadas cartesianas.
\item Calcule $\displaystyle\iiint_D y\,dV$ usando coordenadas cilíndricas.
\end{enumerate}
\end{ejercicio}

\begin{solucion}[Respuesta]
\subsection*{Solución}
\begin{enumerate}[label=\alph*)]
\item En el primer octante ($x\ge0,y\ge0,z\ge0$). El cilindro $x^2+z^2=1$ da $z=\sqrt{1-x^2}$ (rama positiva). El cono $z^2=x^2+y^2$ da $z=\sqrt{x^2+y^2}$. La región está debajo del cilindro y sobre el cono, además $y\ge0$. La proyección sobre $xz$ es un cuarto de círculo $0\le x\le1$, $0\le z\le\sqrt{1-x^2}$. Para cada $(x,z)$, $y$ varía desde $y=0$ hasta $y=\sqrt{z^2-x^2}$ (de la ecuación del cono) siempre que $z\ge x$. Por tanto:
\[
D = \left\{(x,y,z): 0\le x\le1,\; x\le z\le\sqrt{1-x^2},\; 0\le y\le\sqrt{z^2-x^2}\right\}.
\]

\item Usamos coordenadas cilíndricas con $x=r\cos\theta$, $z=r\sin\theta$ (notar que el cilindro es $x^2+z^2=1$). En cilíndricas, el cono $z^2=x^2+y^2$ se convierte en $(r\sin\theta)^2 = (r\cos\theta)^2 + y^2 \Rightarrow r^2\sin^2\theta = r^2\cos^2\theta + y^2 \Rightarrow y^2 = r^2(\sin^2\theta-\cos^2\theta) = -r^2\cos2\theta$. Como $y\ge0$, necesitamos $\cos2\theta\le0$, es decir $\theta\in[\pi/4,\pi/2]$ (pues $0\le\theta\le\pi/2$ en primer octante). Además, $y = \sqrt{-r^2\cos2\theta} = r\sqrt{-\cos2\theta}$. Los límites: $0\le r\le1$, $0\le\theta\le\pi/2$ pero con la restricción adicional de que $z\ge x$ implica $\tan\theta\ge1$, es decir $\theta\ge\pi/4$. También de la desigualdad $y\le\sqrt{z^2-x^2}$ ya está incorporada. El elemento de volumen en cilíndricas (con $x,z$ como polares) es $dV = r\,dr\,d\theta\,dy$ (pues la tercera coordenada es $y$). Entonces:
\[
\iiint_D y\,dV = \int_{\theta=\pi/4}^{\pi/2} \int_{r=0}^{1} \int_{y=0}^{r\sqrt{-\cos2\theta}} y\; (r\,dy\,dr\,d\theta).
\]
Integramos en $y$: $\int_0^{r\sqrt{-\cos2\theta}} y\,dy = \frac12 r^2(-\cos2\theta)$. Luego
\[
\iiint_D y\,dV = \int_{\pi/4}^{\pi/2} \int_0^1 r\cdot \frac12 r^2(-\cos2\theta)\,dr\,d\theta = \frac12 \int_{\pi/4}^{\pi/2} (-\cos2\theta) \int_0^1 r^3 dr\,d\theta.
\]
\[
\int_0^1 r^3 dr = \frac14, \qquad \frac12\cdot\frac14=\frac18.
\]
Entonces
\[
\iiint_D y\,dV = \frac18 \int_{\pi/4}^{\pi/2} (-\cos2\theta)\,d\theta = \frac18 \left[ -\frac12\sin2\theta \right]_{\pi/4}^{\pi/2} = -\frac1{16}\left( \sin\pi - \sin\frac{\pi}{2}\right) = -\frac1{16}(0-1)=\frac1{16}.
\]
Por lo tanto, la integral vale $\boxed{\dfrac{1}{16}}$.
\end{enumerate}
\end{solucion}

\section*{Conclusión}
Las integrales múltiples constituyen una herramienta poderosa para el cálculo de áreas, volúmenes, centros de masa y otras aplicaciones en ciencias e ingeniería. Dominar su cálculo requiere práctica constante y una comprensión geométrica de las regiones de integración. Este material proporciona la teoría esencial y una amplia variedad de ejercicios para alcanzar dicho dominio.

\end{document}